\documentclass{report}

\begin{document}

1. Introduction
Briefly motivate the need for dynamic graph learning.

Identify key challenges: embedding staleness, limited contextual awareness, and poor structural dynamics capture.

Introduce HiCoSTv2 as a solution that integrates multi-scale temporal walks, hierarchical co-occurrence transformers, spectral-temporal ODEs, and robust memory with hard negative mining.

Summarize contributions (novel components, ablation studies, state-of-the-art results).

2. Related Work
Categorize existing methods: discrete-time vs. continuous-time, memory-based (TGN), walk-based (CAWN, CTWalk), transformer-based (DyGFormer), and hybrid approaches (TAWRMAC).

Highlight limitations addressed by HiCoSTv2.

3. Preliminaries
Define temporal graph, temporal walks, anonymization, and problem formulation (link prediction).

4. Methodology (Core of the paper)
4.1 Overview of HiCoSTv2 Architecture
Provide a high-level diagram (flowchart) showing data flow:

Input graph → Multi‑Scale Walk Sampler
Walks → Stabilized Hierarchical Co‑occurrence Transformer (HCT)
Memory → Robust Stability‑Augmented Memory (SAM)
Temporal evolution → Numerically Stabilized ST‑ODE
Fusion → Gated Mutual Refinement Pooling
Training → Hard Negative Mining + Label Smoothing + Cosine LR
4.2 Multi‑Scale Walk Sampler
Describe the generation of short, long, and TAWR walks.

Explain adaptive walk lengths based on node degree/activity.

Introduce temporal bias sampling (Eq. 7 from TAWRMAC) with time noise injection for robustness.

Discuss the masking strategy to handle incomplete walks.

Provide pseudocode or algorithmic description.

4.3 Stabilized Hierarchical Co‑occurrence Transformer (HCT)
Intra‑walk encoding: Transformer with causal masking to respect temporal order.

Co‑occurrence matrix construction using a time‑aware positional kernel.

Inter‑walk transformer with co‑occurrence bias to capture structural correlations.

DropPath regularization for training stability.

Explain how HCT outputs walk‑level embeddings for each node.

4.4 Robust Stability‑Augmented Memory (SAM)
Maintain prototype vectors per node (ablatable).

Query formation combining raw memory, edge features, and time encoding.

Prototype attention to compute candidate memory.

Gated update with residual connection and time‑based decay.

Spectral normalization on prototypes to prevent collapse.

Describe how SAM updates memory after each interaction.

4.5 Numerically Stabilized Spectral‑Temporal ODE (ST‑ODE)
Base ODE dynamics using GRU‑ODE cell.

Spectral regularization via projection onto low‑frequency eigenvectors of the graph Laplacian.

Velocity gating to prevent explosive gradients.

Low‑rank approximation for scalability.

Integration with adjoint method and adaptive step size.

Discuss the update interval strategy to balance efficiency and accuracy.

4.6 Gated Mutual Refinement Pooling
Collect three modalities: SAM memory, HCT structural embeddings, ST‑ODE temporal embeddings.

For each modality, pool walks using attention‑based pooling.

Tri‑modal gating to adaptively fuse modalities per node pair.

Output final source and target embeddings.

4.7 Training Objectives and Optimization
Hard Negative Mining:

Vectorized similarity computation between source nodes and candidate negatives.

Select top‑k hardest negatives based on similarity threshold.

Fallback to random negatives if insufficient.

Label Smoothing to improve generalization.

Cosine learning rate scheduler with warm‑up.

Loss function: Binary cross‑entropy on positive/negative pairs.

4.8 Ablation Design
Clearly define ablation flags:

ablation_sam_prototypes: replace prototypes with single vector.

ablation_hct_hierarchical: disable inter‑walk transformer (simple mean pooling).

ablation_mrp_gating: replace gated fusion with concatenation + MLP.

ablation_walk_multi_scale: use only short walks.

These will be used in Section 5.4 to isolate contribution of each component.

5. Experiments
5.1 Datasets
List datasets (Wikipedia, Reddit, MOOC, LastFM, Enron, etc.) with statistics (nodes, edges, timespan, feature dimensions).

Mention train/validation/test splits (e.g., 70/15/15) and temporal ordering.

5.2 Baselines
Include strong recent methods: TGN, DyGFormer, CAWN, CTWalk, TAWRMAC, etc.

Briefly describe their key mechanisms and why they are relevant.

5.3 Experimental Setup
Implementation details: Framework (PyTorch Lightning), hardware, hyperparameters (learning rate, hidden dim, number of walks, etc.) – can refer to code/config.

Evaluation metrics: AUC‑ROC, Average Precision (AP).

Protocol: Transductive and inductive settings; multiple negative sampling strategies (random, historical, etc.).

5.4 Main Results
Tables comparing HiCoSTv2 against baselines on all datasets for link prediction.

Highlight best performance and relative improvements.

Discuss performance on datasets with/without rich node features.

5.5 Ablation Studies
Run experiments with each ablation flag enabled/disabled.

Present results in a table or bar chart showing contribution of:

Multi‑scale walks

Hierarchical co‑occurrence

Prototype memory

Gated fusion

ST‑ODE

Analyze which component is most critical for different dataset types.

5.6 Analysis of Model Components
Walk analysis: Distribution of walk lengths, restart probabilities learned by TAWR walks.

Memory analysis: Visualize prototype attention weights over time; show how prototypes capture different roles.

Co‑occurrence patterns: Example co‑occurrence matrices for nodes in different communities.

ST‑ODE dynamics: Plot Dirichlet energy over time to show smoothness enforcement.

Hard negative mining: Statistics on mined negatives vs. random negatives.

5.7 Efficiency and Scalability
Compare training/inference time and memory usage with baselines.

Show how vectorized operations and sparse updates improve scalability.

5.8 Hyperparameter Sensitivity
Study impact of key hyperparameters: number of prototypes, number of walks, walk lengths, ODE update interval, etc.

Use line plots to show performance vs. parameter value.

6. Discussion
Summarize findings.

Discuss limitations (e.g., sensitivity to graph size, eigen decomposition cost).

Potential future work: extending to heterogeneous graphs, incorporating edge features more explicitly, or using cheaper spectral approximations.

7. Conclusion
Recap contributions and results.

Conclude with the significance of HiCoSTv2 for dynamic graph learning.


\end{document}